\subsection{固定核 GPPA}
\label{sec:rleb-gppa-main}

\begin{definition}[固定核广义近端点算法]
\label{def:fixed-kernel-GPPA}
固定处处有定义的核 $v:X\to X$。给定允许子关系
$\mathcal J_v\subset J^v_{\lambda F}$，其中
\[
 J^v_{\lambda F}(x)
 =(v+\lambda F)^{-1}(v(x))
 =\{y:v(x)-v(y)\in\lambda F(y)\},
\]
固定核 GPPA 定义为
\begin{equation}
\label{eq:fixed-kernel-GPPA}
 x^{k+1}\in\mathcal J_v(x^k),\qquad
 v(x^k)-v(x^{k+1})\in\lambda F(x^{k+1}).
\end{equation}
取 $\mathcal J_v=J^v_{\lambda F}$ 时，称为完整固定核 GPPA。
\end{definition}

核增量控制本步所选的原图值，但一般不等于物理位移 $x^{k+1}-x^k$。
因而推广普通 PPA 时，需要分别控制核空间中的一步几何和物理空间中的移动。
以下仍以完整真残差 $r_F$ 和物理闭零集 $S$ 为对象，
将普通 PPA 的相应条件替换为固定核条件，并加入物理位移估计。

\subsection{固定核的替换与补充条件}

\begin{assumption}[固定核 GPPA 的 RLEB 条件]
\label{ass:rleb-gppa-supplement}
沿用假设~\ref{ass:rleb-ppa} 的基本对象
$X,F,\lambda,S,U,R,\bar t,\kappa$，以及 $d$、$P_S$、$r_F$ 和 $U_R$ 的约定。
取有限非减函数
\[
 a:[0,R]\to[0,\infty),\qquad
 \omega_v:[0,a(R)]\to[0,\infty),\qquad
 \psi:[0,\bar t]\to[0,\infty),
\]
三者均在零处连续且取值为零。
取有限函数 $h:[0,\infty)^2\to[0,\infty)$，
它对各变量非减，在 $(0,0)$ 连续且 $h(0,0)=0$。定义
\begin{equation}
\label{eq:gppa-physical-moduli}
 \begin{aligned}
 K_v(t)&=\frac{a(t)+\omega_v(a(t))}{2},\\
 H_v(t)&=h(t,K_v(t)),\qquad
 \Theta_v(t)=\sum_{j=0}^{\infty}H_v(\kappa^jt),\quad 0\le t\le R.
 \end{aligned}
\end{equation}
同一个允许关系 $\mathcal J_v$ 满足以下条件。
\begin{enumerate}
\item[(G1)] 对每个 $x\in U_R$，有
$\varnothing\ne\mathcal J_v(x)\subset J^v_{\lambda F}(x)$。
\item[(G2)] 对每个 $x\in U_R$ 和每个 $y\in\mathcal J_v(x)$，
存在 $p\in P_S(x)$，使以下两式同时成立：
\begin{equation}
\label{eq:gppa-same-anchor}
 \begin{aligned}
 \delta&:=\|v(x)-v(p)\|\le a(d(x)),\\
 \langle v(y)-v(p),v(x)-v(y)\rangle
 &\ge\frac{\delta^2-\omega_v(\delta)^2}{4}.
 \end{aligned}
\end{equation}
\item[(G3)] 对上述全部允许转移，有物理位移估计
\begin{equation}
\label{eq:gppa-physical-transfer}
 \|y-x\|\le h\bigl(d(x),\|v(y)-v(x)\|\bigr).
\end{equation}
\item[(G4)] 对每个 $y\in\mathcal J_v(U_R)$，若 $r_F(y)\le\bar t$，则
\begin{equation}
\label{eq:gppa-true-output-EB}
 d(y)\le\psi(r_F(y)).
\end{equation}
此外，输出评价窗满足 $K_v(R)/\lambda\le\bar t$。
\item[(G5)] 对全部 $0<t\le R$，有
\begin{equation}
\label{eq:gppa-compatibility}
 \psi(K_v(t)/\lambda)\le\kappa t,
 \qquad \Theta_v(R)<\infty.
\end{equation}
\end{enumerate}
\end{assumption}

这里 (G1)、(G2)、(G4)、(G5) 分别替换普通 PPA 的
(P1)、(P2)、(P3)、(P4)，而 (G3) 是物理提升的补充条件。
假设~\ref{ass:rleb-gppa-supplement} 独立使用这些替换后的条款，
不要求原关系同时满足替换前的普通 PPA 前件。
特别地，(G3) 必须在全部实际允许转移上成立。
若以核的逆模验证它，估计必须覆盖当前输入与其全部允许输出；
只在 $U\times U$ 上验证逆模，会预先使用尚未证明的输出留域。

\subsection{固定核 GPPA 的局部有限长度}

\begin{theorem}[物理空间中的局部收敛]
\label{thm:rleb-gppa-physical}
在假设~\ref{ass:rleb-gppa-supplement} 下，取 $x^0\in U_R$，记 $d_0=d(x^0)$。
若
\begin{equation}
\label{eq:gppa-physical-budget}
 \Theta_v(d_0)<\dist(x^0,X\setminus U),
\end{equation}
则每条轨道 $x^{k+1}\in\mathcal J_v(x^k)$ 均可无限延拓，
全部点属于 $U_R$，具有物理有限长度，并收敛到某个 $x^\infty\in S\cap U$。
对所有 $k\ge0$，有
\begin{align}
 d(x^k)&\le\kappa^kd_0,\label{eq:gppa-physical-distance}\\
 \|v(x^{k+1})-v(x^k)\|&\le K_v(d(x^k))\le K_v(\kappa^kd_0),\\
 \|x^{k+1}-x^k\|&\le H_v(d(x^k))\le H_v(\kappa^kd_0),\\
 \sum_{k=0}^{\infty}\|x^{k+1}-x^k\|&\le\Theta_v(d_0),\\
 \|x^\infty-x^k\|&\le\Theta_v(d(x^k))\le\Theta_v(\kappa^kd_0).
 \label{eq:gppa-physical-tail}
\end{align}
这些估计统一于全部允许选择，极限点可以随选择改变。
当 $\mathcal J_v=J^v_{\lambda F}$ 时，结论控制完整 GPPA 的全部输出。
\end{theorem}

\begin{proof}
任取 $x\in U_R$ 及 $y\in\mathcal J_v(x)$，并取 (G2) 中的同一个最近零点 $p$。
记 $d=d(x)$ 和 $s_v=\|v(x)-v(y)\|$。
由 (G2) 及柯西不等式，
\[
 s_v^2+\frac{\delta^2-\omega_v(\delta)^2}{4}
 \le\langle v(x)-v(p),v(x)-v(y)\rangle
 \le\delta s_v.
\]
移项得到标量不等式 $(2s_v-\delta)^2\le\omega_v(\delta)^2$，从而
\begin{equation}
\label{eq:gppa-kernel-one-step}
 s_v\le\frac{\delta+\omega_v(\delta)}2\le K_v(d).
\end{equation}
(G1) 给出 $(v(x)-v(y))/\lambda\in F(y)$，故
\[
 r_F(y)\le\frac{s_v}\lambda
 \le\frac{K_v(d)}\lambda\le\frac{K_v(R)}\lambda\le\bar t.
\]
在已验证的输出窗内，(G4)--(G5) 因而给出 $d(y)\le\kappa d$（$d>0$），
而 (G3) 给出 $\|y-x\|\le H_v(d)$。
当 $d=0$ 时，$p=x$，(G2) 强制 $s_v=0$，
(G3) 再强制 $y=x$；距离递推和物理位移界仍然成立。

由 (G1) 可以构造第一步。对任意已构造到 $x^k\in U_R$ 的有限前缀，
每个下一允许输出均满足
\[
 d(x^{k+1})\le\kappa^{k+1}d_0,\qquad
 \|x^{k+1}-x^0\|
 \le\sum_{j=0}^{k}H_v(\kappa^jd_0)
 \le\Theta_v(d_0)<\dist(x^0,X\setminus U).
\]
于是 $x^{k+1}\in U_R$，覆盖条件可以再次使用，归纳给出无限延拓。
在定理~\ref{thm:distance-length} 中取位移包络 $H_v$，
即得轨道在物理空间中的有限长度、点收敛及尾界；同一严格预算保证极限仍在 $U$。
\end{proof}

取 $v=\Id$、$a(t)=t$、$\omega_v=\omega$ 和 $h(t,s)=s$，
便有 $K_v=H_v=B$、$\Theta_v=\Theta$，
假设~\ref{ass:rleb-gppa-supplement} 与定理~\ref{thm:rleb-gppa-physical}
分别恢复普通 PPA 的假设与定理。
一般核的核增量控制真残差，物理提升 (G3) 则把这种控制传递给相邻物理位移。
例如，若在全部允许转移上有
$\|x-y\|\le\eta(\|v(x)-v(y)\|)$，可取 $h(t,s)=\eta(s)$；
非单射核也可通过 $h(t,s)=\eta(s)+\zeta(t)$
控制核所不能区分的方向，只须这些函数满足假设中的条件。


\begin{remark}[非单射核的纤维方向]
物理位移条件不能由核增量的衰减代替。
例如在 $\R^2$ 上取 $F(t,s)=(t,0)$、$v(t,s)=(t,0)$、$\lambda=1$，则
\[
 S=\{0\}\times\R,\qquad
 J^v_F(t,s)=\{(t/2,u):u\in\R\},\qquad r_F(t,s)=|t|.
\]
完整输入覆盖成立，核锚内积等于 $d(x,S)^2/4$。
取 $U=\R^2$、$a(r)=r$、$\omega_v=0$、$\psi(q)=q$、
$R=\bar t=1$、$\kappa=1/2$，
以及 $h(r,q)=r+q$，则 (G1)、(G2)、(G4)、(G5) 均成立。
但 (G3) 不成立：完整关系允许轨道
$x^k=(2^{-k},(-1)^k)$，其核坐标和到 $S$ 的距离均趋于零，
物理点却不收敛。特别地，零输入的完整纤维包含整条 $S$，
而 (G3) 与 $h(0,0)=0$ 会强制零输入的全部允许输出驻定。
\end{remark}
